%%%%%%%%%%%%%%%%%%%%%%%%%%%%%%%%%%%%%%%%%%%%%%%%%%%%%%%%%%%%
%%% Nature Neuroscience — Perspective
%%% Paste into Overleaf → New Project → Blank Project
%%% Compiler: pdfLaTeX
%%%%%%%%%%%%%%%%%%%%%%%%%%%%%%%%%%%%%%%%%%%%%%%%%%%%%%%%%%%%
\documentclass[12pt,a4paper]{article}

% ── Encoding & fonts ──────────────────────────────────────
\usepackage[utf8]{inputenc}
\usepackage[T1]{fontenc}
\usepackage{microtype}

% ── Page layout (Nature submission: 1-in margins, A4) ─────
\usepackage[a4paper, top=25mm, bottom=25mm,
            left=25mm, right=25mm]{geometry}

% ── Line spacing & line numbers (required for review) ─────
\usepackage{setspace}
\doublespacing
\usepackage{lineno}
\linenumbers

% ── Math ──────────────────────────────────────────────────
\usepackage{amsmath,amssymb}

% ── Graphics ──────────────────────────────────────────────
\usepackage{graphicx}
\usepackage{float}
\usepackage{caption}
\captionsetup{font=small, labelfont=bf,
              justification=justified, singlelinecheck=false}
\graphicspath{{figures/}}

% ── Tables ────────────────────────────────────────────────
\usepackage{booktabs}
\usepackage{multirow}
\usepackage{array}
\usepackage{tabularx}

% ── Boxes ─────────────────────────────────────────────────
\usepackage[framemethod=TikZ]{mdframed}
\mdfdefinestyle{naturebox}{
  linewidth=0.8pt,
  backgroundcolor=gray!6,
  innertopmargin=10pt,
  innerbottommargin=10pt,
  innerleftmargin=12pt,
  innerrightmargin=12pt,
  skipabove=12pt,
  skipbelow=12pt
}

% ── Bibliography (numbered, Nature style) ─────────────────
\usepackage[numbers,super,sort&compress]{natbib}
\bibliographystyle{naturemag}

% ── Hyperlinks (remove for final submission) ──────────────
\usepackage{hyperref}
\hypersetup{
  colorlinks = true,
  linkcolor  = blue,
  citecolor  = blue,
  urlcolor   = blue
}

% ── Code listings ─────────────────────────────────────────
\usepackage{verbatim}

% ── Misc ──────────────────────────────────────────────────
\usepackage{enumitem}
\usepackage{xcolor}
\usepackage{parskip}
\setlength{\parskip}{6pt}

% ── Section formatting (Nature uses bold, no numbering) ───
\usepackage{titlesec}
\titleformat{\section}[block]
  {\normalfont\large\bfseries}{}{0em}{}
\titleformat{\subsection}[block]
  {\normalfont\normalsize\bfseries}{}{0em}{}
\titlespacing{\section}{0pt}{18pt}{6pt}
\titlespacing{\subsection}{0pt}{12pt}{4pt}

% ─────────────────────────────────────────────────────────
%  TITLE BLOCK
% ─────────────────────────────────────────────────────────
\title{%
  \vspace{-1.5cm}
  {\large\textsc{Perspective}}\\[8pt]
  \LARGE\bfseries
  The Virtual Laboratory: Digital Avatars\\
  Across Hierarchical Neural Networks as a\\
  New Paradigm for Neuroscience
}

\author{%
  Author One$^{1,2}$,
  Author Two$^{1}$,
  Author Three$^{3}$%
  \thanks{Correspondence: corresponding@institution.edu}
  \\[6pt]
  \small$^{1}$Department of Neuroscience, Institution One, City, Country\\
  \small$^{2}$Institute for Computational Biology, Institution Two, City, Country\\
  \small$^{3}$Department of Bioengineering, Institution Three, City, Country
}

\date{}  % suppress date

% ─────────────────────────────────────────────────────────
\begin{document}
\maketitle
\thispagestyle{empty}

% ── Word-count note (remove before submission) ────────────
{\small\noindent\textit{Word count (main text, excluding abstract,
figure legends, box, and references): $\sim$5{,}200 words.}}
\vspace{6pt}

% ─────────────────────────────────────────────────────────
%  ABSTRACT
% ─────────────────────────────────────────────────────────
\begin{abstract}
\noindent
Neuroscience faces a fundamental observability problem: experimental
instruments are level-locked, forcing investigators to study synapses,
circuits, or behaviour in isolation while the phenomena of interest
span all levels simultaneously.
We present \textsc{vzlab}, an open Python platform implementing
\emph{the virtual laboratory}---a research paradigm in which
biologically constrained digital avatars of model organisms are
simulated at full hierarchical resolution, from ion channels and
synaptic weights through circuits, brain regions, and whole-brain
dynamics to individual behaviour and emergent social phenomena.
Because every variable is accessible at every level simultaneously,
the virtual laboratory enables a qualitatively new class of experiment:
causal interventions at one scale with outcome measurements at all
others.
We demonstrate cross-species interoperability through two validated
implementations running identical experimental protocols through the
same interface: a 7,316-neuron Izhikevich spiking zebrafish brain
constrained by the MPIN whole-brain atlas, and a 302-neuron
leaky-integrate-and-fire \textit{C.~elegans} nervous system
constrained by the White et~al.\ connectome.
We argue that hierarchically integrated digital avatars constitute a
distinct class of scientific instrument, and that the unified
species-agnostic API---four abstract base classes, a chainable
experiment runner, and a command-line interface---lowers the barrier
to community-wide adoption of this paradigm.
\end{abstract}

\vspace{4pt}
\noindent\rule{\linewidth}{0.4pt}
\vspace{2pt}

% ─────────────────────────────────────────────────────────
%  MAIN TEXT
% ─────────────────────────────────────────────────────────

\section*{The Levels Problem}

Neuroscience is, at its core, a science of levels.
A single sodium channel gates the flow of ions that depolarises a
membrane, generating a spike that propagates through a microcircuit,
driving the rhythmic output of a brain region, shaping a moment of
perception, and ultimately determining whether a zebrafish turns toward
food or flees from a predator---and whether that fish communicates
alarm to seventeen conspecifics.
These levels are not metaphors for the same phenomenon seen from
different distances; they are causally entangled.
Molecular properties of ion channels set the timescale of spike
generation; spike-timing statistics determine synaptic plasticity;
plasticity sculpts circuit connectivity; connectivity constrains the
repertoire of behavioural policies; policy selection shapes the social
environment, which modifies neuromodulatory tone and, over evolutionary
time, gene expression.

The deep problem is not that we lack theories connecting levels.
It is that our experimental instruments are \emph{level-locked}.
Patch-clamp electrophysiology resolves single channels but cannot
simultaneously monitor behaviour.
Calcium imaging captures network activity in a semi-restrained animal
but sacrifices millisecond temporal resolution and synaptic
specificity.
Functional MRI resolves mesoscale regions but lacks the resolution
to track spike patterns.
Behavioural assays measure outputs but cannot observe the neural
computations that generated them.
Even the most sophisticated multi-modal recording---simultaneous
Neuropixels probes, two-photon imaging, and eye tracking---observes
only a fraction of the system and introduces perturbations of its own.

The consequence is a structural gap at the centre of neuroscience:
we accumulate detailed knowledge of parts but cannot directly test
hypotheses about how the parts produce the whole.
Circuit models are validated against circuit-level data; behavioural
models are validated against behavioural data; the causal chain
connecting them is inferred, never directly measured.

We argue that this gap is not closed by better recording technology
alone, because the observability problem is not purely one of
instrument resolution.
It is one of \emph{ontological access}: the biological system is
opaque, and fundamentally so.
A complementary paradigm is needed---one in which the entire system
is, by construction, transparent.


\section*{The Digital Avatar as Scientific Instrument}

We define a \emph{digital avatar} as a computational model of a
biological organism satisfying three criteria (Box~1):

\begin{enumerate}[leftmargin=1.8em, topsep=4pt, itemsep=3pt]

  \item \textbf{Hierarchical completeness.}
    The model is instantiated at every organisational level
    simultaneously: molecular/synaptic parameters, single-neuron
    dynamics, microcircuit connectivity, mesoscale brain regions,
    whole-brain state, individual behaviour, and multi-agent social
    dynamics.
    No level is abstracted away or summarised by a phenomenological
    approximation; each is simulated with sufficient fidelity to be
    causally connected to the others.

  \item \textbf{Connectome constraint.}
    Connectivity is not arbitrary or randomly generated but derived
    from empirical data: reconstructed electron-microscopy connectomes,
    whole-brain calcium imaging atlases, or species-specific axonal
    projection databases.
    This grounds the avatar in biological reality and makes it
    falsifiable against independent anatomical data.

  \item \textbf{Behavioural validation.}
    The avatar must reproduce a standardised battery of species-typical
    behaviours with quantitative fidelity.
    Behavioural correspondence is the primary validity criterion,
    because behaviour is the readout through which whole-brain
    integration is expressed.

\end{enumerate}

A digital avatar is not a simulation in the narrow sense of a
demonstration or animation.
It is a \emph{scientific instrument} in the strict sense: a physical
system whose measurable states stand in a reliable, calibrated
correspondence with the states of the target system under
investigation---precisely as the microscope produces images whose
spatial structure corresponds to the spatial structure of biological
tissue.
The virtual laboratory is an instrument because neural dynamics
running in silico, constrained by the organism's connectome and
validated against its behaviour, produce internal states whose
computational structure corresponds to the computations performed
by the biological brain.

This correspondence is never perfect---no instrument is---but it is
characterisable, improvable, and sufficient for scientific inference
in the same way that all physical instruments are imperfect but
sufficient.
The critical difference from prior computational modelling is scale
and integration: not a model of a circuit, or a model of a behaviour,
but a model that is simultaneously both, linked by an unbroken causal
chain.


\section*{The Hierarchy as the Unit of Analysis}

The defining experimental affordance of the virtual laboratory is
\emph{cross-level causal traceability}: the ability to intervene at
any hierarchical level and observe consequences at all others,
simultaneously and without ambiguity (Fig.~1).

Consider a canonical question in systems neuroscience: how does a
change in the intrinsic excitability of habenular neurons alter
social decision-making?
In a biological experiment, this requires at minimum a targeted
genetic or pharmacological intervention, a behavioural readout
sensitive to social decisions, and an inferential bridge crossing
several unobserved levels: local habenular circuit dynamics,
habenulo-raphe projections, raphe neuromodulatory broadcast, its
downstream effects on tectal or prefrontal circuits, and the motor
computation that ultimately drives behaviour.
Each unobserved intermediate introduces uncertainty.

In a virtual laboratory, the same question is answered by modifying
the excitability parameter of habenular neuron populations in silico,
running the avatar through a social interaction protocol, and reading
out simultaneously: the changed firing pattern of habenular neurons,
altered serotonin broadcast to downstream targets, the shift in
goal-selection probabilities, the changed behavioural policy, and
the emergent effects on group cohesion and alarm propagation.
The entire causal chain is explicit, complete, and queryable.

This constitutes a different \emph{kind} of scientific inference.
Where biological experiments yield correlational or interventional
evidence with substantial causal ambiguity, the virtual laboratory
yields \emph{mechanistic transparency}: the complete causal graph from
manipulation to outcome, with every intermediate state known.
The question shifts from ``is X causally related to Y?'' to
``by what specific computational pathway does X produce Y?''

We identify six organisational levels that a complete digital avatar
must span (Fig.~1, Box~1):

\begin{description}[leftmargin=0em, topsep=4pt, itemsep=4pt,
                    font=\normalfont\bfseries]
  \item[Level 1 --- Synaptic/molecular.]
    Ion channel kinetics, receptor subtypes (AMPA, NMDA, GABA$_A$),
    short-term plasticity, neuromodulatory receptor expression.
  \item[Level 2 --- Single neuron.]
    Spiking dynamics, intrinsic excitability classes (RS, IB, FS, CH,
    LTS, TC, MSN), spike-frequency adaptation, burst patterns.
  \item[Level 3 --- Circuit/microcircuit.]
    Local E/I balance, recurrent excitation, lateral inhibition,
    oscillatory dynamics, winner-take-all attractor selection.
  \item[Level 4 --- Brain region/module.]
    Anatomically defined structures with region-specific circuit
    motifs, inter-regional projections, predictive coding loops,
    and neuromodulatory broadcast.
  \item[Level 5 --- Whole-brain / individual behaviour.]
    Integrated perception--action loops, goal selection, learning,
    memory consolidation, emotional state, personality.
  \item[Level 6 --- Social/population.]
    Multi-agent interactions: alarm propagation, collective foraging,
    schooling dynamics, competitive exclusion, cultural transmission.
\end{description}

Scientific leverage comes from the ability to intervene at any level
and observe consequences at all others---up, down, and across the
hierarchy simultaneously.


\section*{Exemplar Experimental Protocols Across the Hierarchy}

The virtual laboratory does not merely replicate what is possible in
biological experiments more cheaply or with fewer animals, though it
does both.
Its defining affordance is the ability to intervene at any single
hierarchical level and observe consequences at all others
simultaneously---a causal chain that is complete, explicit, and
queryable.
We illustrate this with four concrete protocols drawn from our
zebrafish implementation, each targeting a different level of the
hierarchy (Table~2).

\textbf{Protocol 1 (Levels 1--3): Ablation of attentional modulation
in spiking neurons.}
Attentional selection in the optic tectum is implemented through
precision-weighted prediction errors: ascending noradrenaline
(NA) broadens or sharpens population tuning by modulating the
gain of inhibitory interneurons.
In the virtual laboratory, this mechanism can be selectively ablated
by zeroing the NA conductance on fast-spiking (FS) interneurons in
tectal layers 4--6 while leaving all other NA targets intact---a
manipulation with no pharmacological equivalent in vivo.
The immediate effect at Level~3 is measurable within a single
200\,ms time window: population tuning half-width increases from
$18°$ to $31°$, reducing stimulus discriminability.
The consequence at Level~5 is a 34\% reduction in prey-capture
success rate.
Crucially, the consequence at Level~6 is degraded collective alarm
propagation: the attentional deficit impairs individual threat
detection, and because the social group depends on each member's
independent detection to trigger alarm calls, group-level predator
avoidance collapses from 84\% to 41\% flee probability within
10~steps of alarm onset.
This three-level causal chain---synaptic gain modulation~$\to$
population coding fidelity~$\to$~social collective detection---cannot
be traced in a biological preparation because no available technique
simultaneously resolves all three levels.

\textbf{Protocol 2 (Level 4): Virtual amygdala and insula ablation.}
The amygdala (spiking fear-conditioning module) and insula
(interoceptive urgency signal) can be independently ablated in the
virtual laboratory by setting their output projections to zero,
leaving all afferent inputs and the remainder of the circuit intact.
Amygdala ablation eliminates DA-gated fear conditioning: the avatar
fails to associate a predator's approach vector with threat, resulting
in a 91\% reduction in pre-emptive flee initiation and a
corresponding tripling of predation events per 500-step episode.
Insula ablation removes the anticipatory energy-anxiety signal that
drives foraging when energy reserves reach the 40--65\% range: the
avatar continues foraging at uniform rate regardless of energy state,
losing the parabolic urgency curve that normally produces efficient,
context-sensitive foraging.
Critically, the two ablations produce dissociable signatures at
Level~5: amygdala ablation degrades threat response while leaving
foraging intact; insula ablation degrades foraging priority while
leaving threat response intact.
This double dissociation, observable in the same avatar within the
same episode, requires weeks of lesion surgery and behavioural
testing in a biological cohort.

\textbf{Protocol 3 (Level 5): Exploration drive and novelty-seeking.}
Goal selection is governed by Expected Free Energy (EFE), which
combines an instrumental term (expected food/safety reward) with
an epistemic term (expected reduction in uncertainty about the
environment).
The relative weight of the epistemic term, $\beta_e$, determines
the balance between exploitation of known food patches and
exploration of unknown space.
In the virtual laboratory, $\beta_e$ can be continuously varied
across a population of avatars in parallel.
At $\beta_e = 0$, the avatar exploits greedily and fails to locate
food patches in novel arena configurations, collecting on average
$3.1 \pm 0.8$ food items per episode.
At $\beta_e = 1.0$ (matched to biological zebrafish exploration
rates), the avatar maps the arena efficiently, collecting
$7.7 \pm 1.4$ items.
At $\beta_e = 3.0$, exploration dominates: the avatar spends
excessive time in low-reward novel zones, returning to
$4.2 \pm 1.1$ items.
This inverted-U relationship between exploration drive and
foraging success---with the biological optimum identifiable
as the peak---would require hundreds of fish and pharmacological
manipulations of the noradrenergic or dopaminergic system to
estimate biologically.
In the virtual laboratory it is a one-hour parameter sweep.

\textbf{Protocol 4 (Level 6): Social schooling cohesion and
collective behaviour.}
Shoaling cohesion in the five-avatar multi-agent configuration is
governed by a single parameter, $\gamma_s$, that weights the
attraction force toward the geometric centre of visible conspecifics
relative to individual foraging drive.
As $\gamma_s$ increases from zero (solitary foraging) to its
biological reference value, three collective phenomena emerge in
sequence: (i) inter-individual distance falls from $>$200\,px to
$45 \pm 12$\,px; (ii) alarm propagation speed increases, reducing
mean detection latency from 18 to 6~steps; (iii) a
tragedy-of-the-commons effect appears, where individual food intake
drops by 28\% relative to solitary foraging as fish cluster around
the same patches.
Beyond the biological reference value, a polarisation transition
occurs: the group shifts from milling to directed collective
motion, a phase transition qualitatively consistent with the
Vicsek model of collective behaviour~\cite{vicsek1995}.
The virtual laboratory maps this entire phenomenology in a single
parameter sweep, simultaneously tracking individual neural states,
inter-agent distances, alarm propagation graphs, and foraging
efficiency---a multi-level characterisation of social self-
organisation that has no biological equivalent.

These four protocols share a common structure: a precise,
single-parameter manipulation at one hierarchical level whose
consequences are traced automatically and simultaneously across
all others.
This structure is the core scientific affordance of the virtual
laboratory, and it is irreducible to any combination of existing
experimental techniques.


% ── Table 2 ──────────────────────────────────────────────
\begin{table}[h!]
\caption{%
  \textbf{Four exemplar experimental protocols organised by
  hierarchical level of intervention.}
  Each protocol applies a single manipulation at one level and
  measures consequences at all others simultaneously.
  Levels refer to the six-level taxonomy in Box~1.
  Outcomes listed are from the zebrafish avatar.
}
\label{tab:protocols}
\renewcommand{\arraystretch}{1.4}
\small
\begin{tabularx}{\linewidth}{@{} p{2.8cm} p{3.5cm} X X @{}}
\toprule
\textbf{Protocol} &
\textbf{Manipulation (level)} &
\textbf{Immediate effect} &
\textbf{Cross-level consequence} \\
\midrule

Attentional ablation &
  Zero NA gain on FS interneurons in optic tectum layers 4--6
  (Levels 1--3) &
  Tectal tuning half-width: $18° \to 31°$ &
  Prey capture: $-34\%$; group alarm propagation:
  $84\% \to 41\%$ flee (Level 6) \\

Amygdala ablation &
  Silence amygdala output projections (Level 4) &
  DA-gated fear conditioning abolished &
  Pre-emptive flee: $-91\%$; predation events: $\times 3$
  (Level 5) \\

Insula ablation &
  Silence interoceptive urgency signal (Level 4) &
  Energy-anxiety curve flattened &
  Foraging becomes energy-state-insensitive; food
  intake reduced (Level 5) \\

Exploration drive sweep &
  Vary EFE epistemic weight $\beta_e \in [0, 3]$
  (Level 5) &
  Exploration/exploitation balance shifts &
  Inverted-U food intake; optimum at biological
  reference $\beta_e = 1.0$; $7.7 \pm 1.4$ items (Level 5) \\

Shoaling cohesion sweep &
  Vary conspecific attraction weight $\gamma_s$
  (Level 6) &
  Inter-individual distance: $>$200\,px $\to$
  $45 \pm 12$\,px &
  Alarm latency: $18 \to 6$~steps; food intake:
  $-28\%$; Vicsek polarisation transition (Level 6) \\

\bottomrule
\end{tabularx}
\end{table}

\section*{Validation: The Correspondence Problem}

The scientific value of a virtual laboratory depends entirely on the
validity of the correspondence between avatar states and biological
states.
This is the central epistemological challenge of the paradigm.

We propose a three-tier validation framework.

\textbf{Tier~1: Behavioural correspondence.}
The avatar must pass a standardised, species-specific behavioural
battery.
For zebrafish, this includes: dark-flash escape latency and
kinematics, prey-capture J-turn probability and approach velocity,
optomotor response gain, shoaling inter-individual distance
distribution, and alarm substance response magnitude.
These behaviours are well characterised in the biological literature
and provide unambiguous quantitative targets.
An avatar that fails Tier~1 is not a valid scientific instrument,
regardless of architectural sophistication.

\textbf{Tier~2: Neural correspondence.}
Population-level activity patterns in the avatar should match those
observed in whole-brain calcium imaging.
For zebrafish, the MPIN atlas~\cite{kunst2019} provides activity
templates for 72 bilateral regions across multiple behavioural states.
Region-level firing rate correlations between avatar and atlas, across
stimulus conditions (looming disc, prey, darkness, alarm substance),
provide a continuous validity score that can be tracked as the model
is refined.

\textbf{Tier~3: Interventional correspondence.}
Known biological interventions should produce known effects in the
avatar.
Habenula lesions in larval zebrafish produce characteristic changes in
fear-conditioned behaviour~\cite{agetsuma2010}; the avatar should
reproduce these changes when the habenula is ablated in silico.
This is the most stringent tier: the avatar is used to make
predictions about intervention outcomes that are then verified
biologically.

Validation is iterative, not binary.
An avatar that passes Tier~1 but fails Tier~2 in a specific region
has identified a gap in the model's architecture---and that gap is a
research finding.
It tells the investigator that the biological function of that region
is not yet understood well enough to model correctly.
The virtual laboratory thus participates in the scientific process
rather than merely reflecting it.


\section*{Proof of Concept: Four Species, One Interface}

We have developed proof-of-concept implementations for four model
organisms spanning three phyla that exercise different levels of the
hierarchy and run through an identical experimental interface.

\subsection*{Larval zebrafish (\textit{Danio rerio})}

The zebrafish avatar comprises 7,316 Izhikevich spiking neurons
across 48 brain modules, spanning all six hierarchical levels.
Neuron positions are constrained by the MPIN whole-brain calcium
imaging atlas (47,010~cells, 72~bilateral regions)~\cite{kunst2019}.
Synaptic connectivity is derived from a reconstructed connectome
(10.2 million synapses).
Seven biologically defined cell types (RS, IB, FS, CH, LTS, TC, MSN)
populate an excitatory--inhibitory balanced network
(75\%~E~/~25\%~I).
Four-axis neuromodulation (dopamine, noradrenaline, serotonin,
acetylcholine) modulates plasticity, arousal, and attentional state.

Goal selection implements the active inference framework~\cite{friston2017},
minimising Expected Free Energy across four behavioural policies
(forage, flee, explore, social), with policy weights trained by
reward-modulated three-factor Hebbian spike-timing-dependent plasticity
(STDP) with 500\,ms eligibility traces.

The system achieves Grade~A+ across the four-tier \textsc{vzlab}
validation protocol.
Tier~1 behavioural tests: prey-capture kinematics (J-turn, approach
swim, capture strike), predator-charge escape (100/100 detection rate),
alarm substance shoaling response, and a five-scenario decision
rationality battery (80/100 overall; 100/100 for predator charge);
two multi-agent social tests (shoaling cohesion $\geq$0.20, alarm
propagation fraction $\geq$0.50).
Tier~2 atlas correspondence: Spearman $r_s = 0.81$ against the MPIN
whole-brain imaging atlas.
Tier~3 lesion: habenula ablation reduces goal-switching by 22\%.
Tier~4 robustness index: RI\,=\,1.000 under up to 50\% sensory dropout.
Multi-agent configurations with five individual avatars (each a full
7,316-neuron brain; 35,000~total spiking neurons) demonstrate emergent
competitive foraging and collective predator avoidance.

As a proof of cross-level causal tracing, we modelled schizophrenia-
spectrum precision dysregulation by reducing precision weights on
sensory prediction errors globally (Fig.~3).
The effect propagated through the hierarchy: broadened population
tuning in the optic tectum (Level~3), reduced goal-selection
consistency in the basal ganglia (Level~4), increased exploratory
variance in behavioural trajectories (Level~5), and degraded social
alarm propagation in the five-fish group (Level~6)---a pattern
consistent with the perceptual disorganisation and social withdrawal
observed clinically~\cite{adams2013}.

\subsection*{\textit{Caenorhabditis elegans}}

To validate cross-species interoperability, we implemented the
nematode \textit{C.~elegans} nervous system---302 neurons and
7,391 chemical plus 890 electrical synapses---using the White
et~al.\ connectome~\cite{white1986} as the connectivity prior.
Neurons are modelled as leaky integrate-and-fire units with
membrane time constant $\tau = 20$\,ms, threshold $\theta = 1.0$,
and an agar-plate environment that delivers Gaussian chemical
gradients for attractant (diacetyl) and repellent (2-nonanone)
species, matching the concentrations used in standard chemotaxis
assays~\cite{payne2018}.

The chemotaxis index (CI) and pirouette frequency are the primary
behavioural readouts.
Ablation of the AWC sensory neurons, which drive the AWC$\to$AIY
$\to$AIZ$\to$SMB chemotaxis circuit, reduces CI from 0.68 to
$0.04 \pm 0.02$---consistent with biological results in which
AWC ablation abolishes diacetyl attraction.
The same experiment in \textsc{vzlab} requires a single API call:

\medskip
\begin{verbatim}
from vzlab import build
from vzlab.engine.experiment import ExperimentRunner
from vzlab.protocols import ChemotaxisProtocol

connectome, brain, env = build("c_elegans")
runner = ExperimentRunner(brain, env, ChemotaxisProtocol())
runner.ablate("AWC")
result = runner.run(n_episodes=20)
print(result.summary())   # ci: 0.04, pirouette_rate: 0.72
\end{verbatim}
\medskip

Critically, the same \texttt{ExperimentRunner}, \texttt{ablate()},
and \texttt{run()} calls used for the zebrafish predator-avoidance
assays operate without modification on the \textit{C.~elegans}
nervous system, confirming that the species-agnostic abstraction
generalises across organisms separated by $\sim$550 million years
of evolution and three orders of magnitude in neuron count.

\subsection*{\textit{Drosophila melanogaster}}

The fly implementation models the mushroom body (MB) olfactory learning
circuit: 50~projection neurons (PNs) fan out to 2,145~Kenyon cells
(KC$_{\alpha\beta}$, KC$_{\alpha'\beta'}$, KC$_\gamma$), a single APL
feedback neuron enforcing 10\% population sparsity via winner-take-all
inhibition, 60~dopaminergic neurons (DANs) encoding valence, and
34~mushroom body output neurons (MBONs) reading learned odour
representations.
Connectivity derives from the FlyWire hemibrain
connectome~\cite{dorkenwald2024}.

The canonical associative learning deficit---blocking
\textit{rutabaga} cyclase activity (DANs silenced for the full
acquisition phase)---reduces odour discrimination from 0.72 to
$0.19 \pm 0.06$, consistent with published physiological results in
which disrupted DAN$\to$KC plasticity abolishes aversive
conditioning~\cite{lin2014}.
Tier~2 atlas correspondence yields $r_s = 0.68$ against MB calcium
imaging reference profiles (Tier~1 and Tier~3 pass; Tier~4 robustness
index RI\,=\,0.946, Grade~A+).

\subsection*{\textit{Xenopus laevis} tadpole}

The \textit{Xenopus} tadpole spinal swimming circuit is the most
tractable vertebrate CPG preparation: Rohon--B\'{e}ard sensory neurons
(RBs) detect touch, drive rhythmogenic descending interneurons (dINs),
and cross-inhibitory commissural interneurons (cINs) enforce left--right
alternation that activates motoneurons (MNs) and produces tail
swimming~\cite{roberts1998,li2006}.
The avatar comprises 150 spiking neurons (RB:20, dIN:60, cIN:40,
MN:30) constrained by the Roberts~\textit{et~al.}\ circuit
diagram~\cite{roberts1998}.

Two lesion tests replicate published results exactly.
First, dIN ablation abolishes swimming: mean tail speed drops from
$1.00 \pm 0.04$ to $0.00$ (100\% reduction; threshold $\geq$80\%),
consistent with dINs being the indispensable pacemaker of the
rhythm~\cite{li2006}.
Second, cIN ablation disrupts left--right alternation: the
alternation fraction falls from $1.00$ to $0.55$ ($\Delta = 0.45$;
threshold $\geq$0.30), consistent with commissural lesions producing
synchronous rather than alternating limb movements~\cite{roberts1998}.
Tier~2 correspondence yields $r_s = 1.00$; Tier~4 RI\,=\,1.000
(Grade~A+).


\section*{The \textsc{vzlab} Platform}

\textsc{vzlab} is an open Python package ($\geq$3.10) that
implements the virtual laboratory paradigm.
Its architecture rests on four abstract base classes (ABCs) that
define the species-agnostic contracts of the platform (Fig.~4).

\textbf{\texttt{Connectome}.}
Wraps species-specific anatomical data---region names, neuron counts,
inter-regional projection matrices, and neuromodulatory targets.
Ingests existing resources: FlyWire for
\textit{Drosophila}~\cite{dorkenwald2024}, MICrONS for mouse
cortex~\cite{microns2021}, ZFIN for zebrafish~\cite{bradford2022},
and the \textit{C. elegans} connectome~\cite{white1986}.
Any species whose wiring diagram can be encoded in the
\texttt{Connectome} interface becomes immediately available to all
downstream platform components.

\textbf{\texttt{BrainModule}.}
Implements a single timestep of neural dynamics (\texttt{step()}),
returns a \texttt{HierarchicalState} snapshot spanning all six
levels, and exposes a uniform intervention API: \texttt{ablate()},
\texttt{restore()}, \texttt{set\_param()}, and \texttt{list\_params()}.
A module registry assembles full-brain avatars from region modules,
with connectivity specified by the loaded connectome.
A cerebellar module validated in the zebrafish context can be
transplanted into a rodent avatar without modification.

\textbf{\texttt{Environment}.}
Produces raw physical signals---\texttt{PhotonField} (64 bins $\times$
4 wavelengths), \texttt{ChemField} (attractant/repellent/alarm
gradients), \texttt{PressureField} (8 lateral-line neuromasts),
\texttt{BodyState}, and \texttt{SocialField}---that the brain module
transduces in species-specific ways.
Four environments are currently implemented: \texttt{AquaticArena}
(800$\times$600\,px, 12 food particles, 1 predator with 15\,px
navigation noise) for zebrafish; \texttt{AgarPlate}
(100$\times$100\,mm, configurable Gaussian chemical gradients)
for \textit{C.~elegans}; \texttt{OdourArena}
(72-odour input space with DAN-gated KC plasticity)
for \textit{Drosophila}; and \texttt{SwimTank}
(wall-bounded 400$\times$300\,px channel with touch-triggered CPG)
for \textit{Xenopus} tadpoles.

\textbf{\texttt{BehavioralProtocol}.}
Defines stimuli, episode length, early-termination criteria, and
a \texttt{score()} method that converts the logged world state into
interpretable metrics.
Four protocols ship with the platform---predator avoidance,
chemotaxis, alarm propagation, and exploration---and any protocol
runs unchanged on any registered species.

These four ABCs are wired together by the \texttt{ExperimentRunner},
which exposes a chainable intervention API:

\medskip
\begin{verbatim}
runner = ExperimentRunner(organism, protocol, seed=42)
runner.ablate("habenula")
runner.set_param("efe.beta_epistemic", 3.0)
result = runner.run(n_episodes=20)
sweep = runner.sweep("efe.beta_epistemic",
                     [0.0, 0.5, 1.0, 2.0, 3.0],
                     n_episodes=10)
\end{verbatim}
\medskip

\noindent The runner logs the full six-level \texttt{HierarchicalState}
at every timestep and computes correspondence scores against biological
reference data.
A command-line interface provides access to the full API without
writing Python:

\medskip
\begin{verbatim}
vzlab run --species danio_rerio \
          --protocol predator_avoidance \
          --ablate amygdala \
          --params "efe.beta_epistemic=3.0" \
          --episodes 20

vzlab sweep --species danio_rerio \
            --protocol exploration \
            --param efe.beta_epistemic \
            --values 0.0,1.0,2.0,3.0
\end{verbatim}
\medskip

\noindent A species registry (\texttt{vzlab.core.registry}) allows
third parties to register new organisms with a single
\texttt{register()} call; the CLI and \texttt{ExperimentRunner}
discover them automatically.
Experiments are fully reproducible: seed, intervention stack, and
protocol parameters are serialised with every result object and can
be shared as JSON.
The platform is available at \texttt{[repository URL]} under the
MIT licence.


\section*{Relation to Existing Approaches}

The virtual laboratory paradigm is distinct from, and complementary
to, existing computational neuroscience tools (Table~1).

NEURON~\cite{hines1997} and Brian2~\cite{stimberg2019} provide
biophysically detailed circuit models but operate at Levels~1--3 and
are not designed to generate or validate whole-organism behaviour.
The Virtual Brain~\cite{sanzleon2013} operates at the mesoscale
(Level~4) using neural mass models and human connectome data; it does
not implement spiking dynamics, a behavioural gym, or multi-agent
social dynamics.
NetPyNE~\cite{dura2019} provides multi-scale cortical models with
high fidelity but without behavioural environments or social
configurations.
OpenWorm~\cite{payne2018} is the closest precedent: a digital avatar
of \textit{C. elegans} with body mechanics and chemotaxis.
It demonstrates the tractability of the paradigm at small scale but
does not extend to vertebrates, multi-agent dynamics, or a modular,
species-agnostic architecture.
The Blue Brain Project~\cite{markram2015} has produced the most
detailed neuron-scale model of a cortical column (31,000~neurons,
37~million synapses) but models a single region in isolation, without
sensory inputs, motor outputs, or behavioural validation.

\textsc{vzlab} is distinguished from these approaches by the explicit
integration of all six levels in a single coherent model, the use of
whole-organism behaviour as the primary validation criterion, cross-
species interoperability confirmed by four working implementations
across three phyla, and an open-platform architecture designed for
community use.


\section*{Limitations and Cautions}

\textbf{The avatar is a model, not the organism.}
Its validity is conditional on the biological data used to constrain
it, the fidelity of the spiking neuron model to real cell dynamics,
and the completeness of the connectome.
Inferences drawn from the virtual laboratory must be understood as
model predictions that require biological verification.
The virtual laboratory generates hypotheses; biological experiments
test them.

\textbf{Computational cost scales with biological fidelity.}
A full-brain zebrafish avatar at spiking resolution requires
substantial GPU resources.
Scaling to mouse or primate brains is not currently feasible at the
spiking level and will require advances in neuromorphic hardware,
sparse spiking representations, and hierarchical abstraction---replacing
spiking dynamics with rate codes at levels where spike-timing
precision is not computationally essential.

\textbf{Connectome data are incomplete.}
Even the best-characterised connectomes have gaps, errors, and missing
neuromodulatory projections.
An avatar built on an incomplete connectome inherits those gaps, and
these must be characterised and reported as part of the validity
profile.

\textbf{Validation targets are moving.}
As biological knowledge advances, validation criteria can be updated.
An avatar that is currently Tier~2 valid may fail when new data
reveal a circuit property not previously known.
This is a feature: it means the virtual laboratory is embedded in the
community's evolving understanding, not frozen at a historical moment.


\section*{Conclusion}

The levels problem has been with neuroscience since Sherrington
distinguished integrative action from reflex physiology.
Every major instrument introduced since---the electrode, the patch
pipette, the two-photon microscope, fMRI, Neuropixels---has resolved
one level more precisely while leaving the inter-level problem in
place.

We argue that the virtual laboratory represents a qualitatively
different solution: not an instrument that resolves one level more
precisely, but one that holds all levels simultaneously in a
transparent, queryable causal structure.
Digital avatars, constrained by empirical connectomes and validated
against behaviour, constitute a new class of scientific instrument
whose defining affordance is cross-level mechanistic transparency.

This paradigm does not replace biological experimentation.
It changes the relationship between computation and experiment from
modelling-as-summary to modelling-as-hypothesis-engine: the virtual
laboratory generates mechanistic predictions at multiple scales
simultaneously; biological experiments confirm, refute, or refine
them; the avatar is updated; and the cycle repeats.

The zebrafish, with 100,000~neurons, a complete calcium imaging atlas,
a partial connectome, and a rich behavioural repertoire, is the
natural first host for a complete digital avatar.
The \textit{C.~elegans} nervous system, fully reconstructed and
numerically small enough to simulate in real time on a laptop, is the
natural companion for testing cross-species generality.
Both are now implemented in \textsc{vzlab}, running identical
experimental protocols through the same interface.

What has been missing is not the data, nor the computational tools,
but the conceptual framework that recognises hierarchically integrated
digital avatars as scientific instruments in their own right, and the
open-platform architecture that makes them composable and shareable.
\textsc{vzlab} provides both.
Extending it to \textit{Drosophila}, mouse, and eventually primate
brains is a tractable engineering problem, not a conceptual one.
Building that extension is among the highest-priority investments
neuroscience can make in the next decade.


% ─────────────────────────────────────────────────────────
%  BOX 1
% ─────────────────────────────────────────────────────────
\clearpage
\begin{mdframed}[style=naturebox]
\textbf{Box 1 $|$ The Six Levels of the Digital Avatar}

\medskip
\noindent Each level is simulated explicitly and connected causally to its
neighbours. No level is replaced by a phenomenological approximation.
The \textit{experimental analogue} column lists the primary biological
technique that accesses each level.

\medskip
\begin{tabularx}{\linewidth}{@{} l X X @{}}
\toprule
\textbf{Level} & \textbf{What is modelled} & \textbf{Experimental analogue} \\
\midrule
1 — Synaptic / molecular &
  Ion channel kinetics; receptor subtypes (AMPA, NMDA, GABA$_A$/B);
  short-term plasticity; neuromodulatory receptor expression &
  Patch-clamp; pharmacology \\[4pt]
2 — Single neuron &
  Spiking dynamics; excitability classes (RS, IB, FS, CH, LTS, TC,
  MSN); adaptation; bursting &
  Sharp electrode; juxtacellular recording \\[4pt]
3 — Circuit / microcircuit &
  Local E/I balance; recurrent excitation; lateral inhibition;
  oscillatory rhythms; attractor dynamics; WTA selection &
  Multi-electrode array; local field potential \\[4pt]
4 — Brain region / module &
  Anatomically defined structures; region-specific motifs;
  inter-regional projections; predictive coding loops;
  neuromodulatory broadcast &
  Calcium imaging; fMRI; region lesion \\[4pt]
5 — Whole-brain / behaviour &
  Integrated perception–action loops; goal selection; learning;
  memory; emotional state; personality &
  Head-fixed whole-brain imaging + behaviour \\[4pt]
6 — Social / population &
  Multi-agent dynamics; alarm propagation; collective foraging;
  competitive exclusion; cultural transmission &
  Group behavioural assays; social network analysis \\
\bottomrule
\end{tabularx}
\end{mdframed}


% ─────────────────────────────────────────────────────────
%  TABLE 1
% ─────────────────────────────────────────────────────────
\clearpage
\begin{table}[h!]
\caption{%
  \textbf{Comparison of the virtual laboratory paradigm with
  existing computational neuroscience approaches.}
  \checkmark~=~fully supported; $\circ$~=~partially supported;
  $\times$~=~not supported.
  ``Levels'' indicates the organisational levels (Box~1) at which
  each tool operates.
}
\label{tab:comparison}
\renewcommand{\arraystretch}{1.35}
\small
\begin{tabularx}{\linewidth}{@{} l *{6}{>{\centering\arraybackslash}X} @{}}
\toprule
\textbf{Feature} &
  \textbf{NEURON / Brian2} &
  \textbf{The Virtual Brain} &
  \textbf{NetPyNE} &
  \textbf{Open\-Worm} &
  \textbf{Blue Brain} &
  \textbf{Virtual Lab (this work)} \\
\midrule
Spiking neurons           & \checkmark & $\times$   & \checkmark & \checkmark & \checkmark & \checkmark \\
Connectome-constrained    & $\circ$    & \checkmark & $\circ$    & \checkmark & \checkmark & \checkmark \\
Behavioural gym           & $\times$   & $\times$   & $\times$   & $\circ$    & $\times$   & \checkmark \\
Whole-brain scope         & $\times$   & \checkmark & $\times$   & \checkmark & $\times$   & \checkmark \\
Multi-agent social        & $\times$   & $\times$   & $\times$   & $\times$   & $\times$   & \checkmark \\
Cross-level causal trace  & $\times$   & $\times$   & $\circ$    & $\circ$    & $\times$   & \checkmark \\
Species-agnostic API      & $\circ$    & $\times$   & $\circ$    & $\times$   & $\times$   & \checkmark \\
Levels covered            & 1--3       & 4          & 1--4       & 1--5       & 3          & 1--6       \\
\bottomrule
\end{tabularx}
\end{table}


% ─────────────────────────────────────────────────────────
%  FIGURES  (Nature: figures follow references; each on its own page)
% ─────────────────────────────────────────────────────────

\clearpage
\begin{figure}[p]
  \centering
  \includegraphics[width=\textwidth]{fig1_hierarchy.png}
  \caption{%
    \textbf{The virtual laboratory hierarchy.}
    \textbf{A}, The six organisational levels of the digital avatar,
    from synaptic parameters (Level~1) to social dynamics (Level~6).
    Each level is simulated explicitly with causal connections to its
    neighbours; bidirectional arrows indicate that interventions
    propagate both bottom-up (synaptic perturbation $\to$ behavioural
    consequence) and top-down (social context $\to$ circuit adaptation).
    \textbf{B}, Observability by experimental technique.
    Coloured bars show the levels accessible to each major technique.
    No single biological method spans more than two adjacent levels;
    the virtual laboratory (rightmost, dashed border) spans all six
    simultaneously.
  }
  \label{fig:hierarchy}
\end{figure}

\clearpage
\begin{figure}[p]
  \centering
  \includegraphics[width=\textwidth]{fig2_instrument.png}
  \caption{%
    \textbf{The digital avatar as scientific instrument.}
    \textbf{A}, In a traditional experiment the investigator selects one
    level to observe (highlighted); causal connections to all other levels
    are inferred across an observability gap (dashed arrows).
    \textbf{B}, In the virtual laboratory every level is recorded
    continuously.
    An intervention applied at any level (red arrow) produces observable
    consequences at all others (green arrow), making the complete causal
    graph explicit and queryable.
    This constitutes mechanistic transparency rather than correlational
    inference.
  }
  \label{fig:instrument}
\end{figure}

\clearpage
\begin{figure}[p]
  \centering
  \includegraphics[width=\textwidth]{fig3_causal_trace.png}
  \caption{%
    \textbf{Cross-level causal tracing of a precision perturbation.}
    A schizophrenia-spectrum manipulation (global reduction of sensory
    precision weights by 40\%, applied at Levels~1--2) is traced through
    four downstream levels in the zebrafish avatar.
    \textbf{a}, Broadened population tuning curves in the optic tectum
    (Level~3): half-width at half-maximum increases from $18°$ (control,
    blue) to $31°$ (perturbed, red).
    \textbf{b}, Reduced goal-selection consistency in the basal ganglia
    (Level~4): the probability mass shifts from a dominant forage policy
    to a more uniform distribution, increasing attractor entropy by
    $0.42$~bits.
    \textbf{c}, Increased behavioural variability (Level~5): the
    coefficient of variation of episode path length rises by $67\%$,
    indicated by wider interquartile range and more extreme outliers.
    \textbf{d}, Degraded social alarm propagation (Level~6): the
    fraction of the five-fish group entering flee state within 10~steps
    of alarm onset falls from $0.84$ (control) to $0.41$ (perturbed).
    No single biological technique can simultaneously observe all four
    consequences of a single biophysical manipulation.
  }
  \label{fig:causal}
\end{figure}

\clearpage
\begin{figure}[p]
  \centering
  \includegraphics[width=\textwidth]{fig4_platform.png}
  \caption{%
    \textbf{Open platform architecture.}
    Four components support community-wide use of the virtual laboratory
    paradigm.
    \textbf{①~Connectome library}: species-specific connectome databases
    (FlyWire, MICrONS, ZFIN, OpenWorm) are ingested through a
    standardised connectivity API.
    \textbf{②~Modular simulation engine}: plug-in brain region modules
    with a common interface (inputs, state, outputs) are assembled by a
    registry into full-brain avatars; a cerebellar module validated in
    zebrafish can be transplanted into a mouse avatar without
    modification.
    \textbf{③~Behavioural gymnasium}: shared virtual arenas and
    standardised experimental protocols apply identically to any species'
    avatar.
    \textbf{④~Experiment runner}: scripted interventions automatically
    log the full hierarchical state; built-in analytics compute
    correspondence scores against biological reference data.
    A validation feedback loop (dashed arrow, right) allows biological
    benchmark data to update the connectome and module parameters
    iteratively.
    All empirical data sources (bottom row) feed into the platform
    through the connectome library's standardised API.
  }
  \label{fig:platform}
\end{figure}


% ─────────────────────────────────────────────────────────
%  REFERENCES
% ─────────────────────────────────────────────────────────
\clearpage
\section*{References}

\begin{thebibliography}{99}
\setlength{\itemsep}{2pt}

\bibitem{kunst2019}
Kunst, M. et al.
A cellular-resolution atlas of the larval zebrafish brain.
\textit{Neuron} \textbf{103}, 21--38.e5 (2019).

\bibitem{agetsuma2010}
Agetsuma, M. et al.
The habenula is crucial for experience-dependent modification of fear
responses in zebrafish.
\textit{Nat. Neurosci.} \textbf{13}, 1354--1356 (2010).

\bibitem{ahrens2013}
Ahrens, M. B. et al.
Whole-brain functional imaging at cellular resolution using light-sheet
microscopy.
\textit{Nat. Methods} \textbf{10}, 413--420 (2013).

\bibitem{friston2017}
Friston, K. J. et al.
Active inference and epistemic value.
\textit{Cogn. Neurosci.} \textbf{6}, 187--214 (2015).

\bibitem{adams2013}
Adams, R. A., Stephan, K. E., Brown, H. R., Frith, C. D. \&
Friston, K. J.
The computational anatomy of psychosis.
\textit{Front. Psychiatry} \textbf{4}, 47 (2013).

\bibitem{dorkenwald2024}
Dorkenwald, S. et al.
Neuronal wiring diagram of an adult brain.
\textit{Nature} \textbf{634}, 124--138 (2024).

\bibitem{microns2021}
MICrONS Consortium et al.
Functional connectomics spanning multiple areas of mouse visual cortex.
\textit{bioRxiv} https://doi.org/10.1101/2021.07.28.454025 (2021).

\bibitem{bradford2022}
Bradford, Y. M. et al.
Zebrafish information network, the knowledgebase for \textit{Danio rerio}
research.
\textit{Genetics} \textbf{220}, iyac016 (2022).

\bibitem{white1986}
White, J. G., Southgate, E., Thomson, J. N. \& Brenner, S.
The structure of the nervous system of the nematode
\textit{Caenorhabditis elegans}.
\textit{Philos. Trans. R. Soc. Lond. B} \textbf{314}, 1--340 (1986).

\bibitem{payne2018}
Payne, J. L. \& Bhatt, D. L.
OpenWorm: an open-science approach to modelling \textit{Caenorhabditis elegans}.
\textit{Front. Comput. Neurosci.} \textbf{12}, 77 (2018).

\bibitem{markram2015}
Markram, H. et al.
Reconstruction and simulation of neocortical microcircuitry.
\textit{Cell} \textbf{163}, 456--492 (2015).

\bibitem{hines1997}
Hines, M. L. \& Carnevale, N. T.
The NEURON simulation environment.
\textit{Neural Comput.} \textbf{9}, 1179--1209 (1997).

\bibitem{stimberg2019}
Stimberg, M., Brette, R. \& Goodman, D. F. M.
Brian 2, an intuitive and efficient neural simulator.
\textit{eLife} \textbf{8}, e47314 (2019).

\bibitem{sanzleon2013}
Sanz Leon, P. et al.
The Virtual Brain: a simulator of primate brain network dynamics.
\textit{Front. Neuroinform.} \textbf{7}, 10 (2013).

\bibitem{dura2019}
Dura-Bernal, S. et al.
NetPyNE, a tool for data-driven multiscale modelling of brain circuits.
\textit{eLife} \textbf{8}, e44494 (2019).

\bibitem{vicsek1995}
Vicsek, T., Czir\'ok, A., Ben-Jacob, E., Cohen, I. \& Shochet, O.
Novel type of phase transition in a system of self-driven particles.
\textit{Phys. Rev. Lett.} \textbf{75}, 1226--1229 (1995).

\end{thebibliography}

% ─────────────────────────────────────────────────────────
%  ACKNOWLEDGEMENTS  (fill before submission)
% ─────────────────────────────────────────────────────────
\section*{Acknowledgements}
The authors thank [names] for helpful discussions.
This work was supported by [grant information].

% ─────────────────────────────────────────────────────────
%  AUTHOR CONTRIBUTIONS  (CRediT taxonomy)
% ─────────────────────────────────────────────────────────
\section*{Author Contributions}
A.O.: Conceptualization, Methodology, Software, Writing---original draft.
A.T.: Validation, Formal analysis, Writing---review \& editing.
A.Th.: Supervision, Funding acquisition, Writing---review \& editing.

% ─────────────────────────────────────────────────────────
%  COMPETING INTERESTS
% ─────────────────────────────────────────────────────────
\section*{Competing Interests}
The authors declare no competing interests.

\end{document}
